\section{本讲小结}

薄透镜把两次近轴折射合并到同一量距平面，要求厚度修正相对所求精度可略。焦距由形状与相对折射率共同决定；两侧介质不同，应使用 $n_u/u+n_v/v=\Phi$，分别定义物方、像方焦距，横向放大率也保留介质因子。

组合系统沿实际传播顺序处理。前一级的会聚目标在后一级之后时，它是后一级的虚物；镀银系统还要在反射后重新指定传播方向，并计入回程折射。主平面提供等效量距基准，首末顶点到焦点的距离与有效焦距一般不同。用矩阵计算时，应补齐物、像两端传播，对完整矩阵施加成像条件。

紧贴光焦度直接相加，有限间隔还包含传播引起的高度变化。无焦系统没有有限焦距，仍可改变束宽和倾角。光斑反求成像保留屏上有向高度，先区分交叉与未交叉的几何支路；轴向位移、物像距离的导数和横向放大率也分别按其定义计算。

消色差同时满足总光焦度与一阶色散两项约束，不能把正负焦距的大小简单相加。特殊齐明位置依赖精确距离比，精确面形依赖等光程及传播支路；球面只在局部近轴曲率上与精确曲面一致。检查图中的实际箭头、延长线、孔径和外部焦点条件，才能把代数解落实为允许的光路。
